\documentclass[../../../include/open-logic-section]{subfiles}

\begin{document}
\olfileid{sth}{ordinals}{zfm}
\olsection{$\ZFminus$۔ اک اہم پڑاؤ}

بدل رکھن دے اصول دا جواز کیویں دینا اے (جے دے وی سکدے آں)، ایہ سدھا
سوال نہیں۔ سو ایہ گل اسیں \olref[replacement][]{chap} لئی رکھاں گے۔
پر بدل رکھن دا اصول شامل کر کے اسیں اک ہور اہم پڑاؤ اُتے پہنچ گئے آں۔
ہُن ساڈے کول نظریے $\ZFminus$ لئی لوڑیندے سارے مسلّمے نیں۔ تفصیل ایہ اے:

\begin{defn}
نظریے $\ZFminus$ دے مسلّمے ایہ نیں۔ عنصراں نال تعیّن، اتحاد، جوڑیاں،
ذیلی سیٹاں دے سیٹ، لامتناہیت، تے وکھ کرن تے بدل رکھن دے خاکیاں دیاں
ساریاں صورتاں۔ دوجے لفظاں وچ، $\ZFminus$، $\Zminus$ وچ بدل رکھن
دا اصول شامل کردا اے۔
\end{defn}
\noindent
ایہ ناں \emph{زرمَیلو--فرینکل} دے سیٹ نظریے لئی اے (\emph{گھٹاؤ}
اوہ شے جیہڑی اسیں اگے ویکھاں گے)۔ فرینکل ایس عزت دا حق دار اے، کیوں
جے \citeyear{Fraenkel1922} وچ بدل رکھن دا اصول بیان کرن دا سہرا اوہدے
سر اے، بھانویں پہلا ٹھیک ٹھیک بیان \citet{Skolem1922} نے دتا سی۔

\end{document}
